% Zwei Tabellen mit mehrzeiliger Legende auf DERSELBEN Seite: jeder
% Legendenkopf ("Tabelle N") traegt sein festes "\\[1pt]" zum Legendentext,
% die Seite enthaelt den Sollwert+1pt-Abstand also ZWEIMAL. Regressionsnetz
% fuer A13: bis zum Fix galten die beiden gleichen Luecken als Cluster und
% damit als Kennzahlenverstoss (gefunden an den tabellenreichen
% Produktionsdokumenten des KIVAT-Forks, 2026-08-05).
\documentclass{swisstex}
\setrunningtitle{Fixture Legendenpaar}
\begin{document}
\section{Abschnitt}
Grundtext auf dem Raster, damit die Seite regulaer beginnt und die
Marginalspalte neben beiden Tabellen bestueckt wird.

\begin{swisstable}{Erste Legende, lang genug fuer mehrere Zeilen in der Marginalspalte}
\begin{tabularx}{\textcolumn}[t]{@{}S{30mm}@{\hspace{4mm}}L@{}}
\toprule
\tabhead{Merkmal} & \tabhead{Beschreibung} \\
\midrule
Erstes & Ein Wert im Tabellenkoerper \\
Zweites & Noch ein Wert \\
\bottomrule
\end{tabularx}
\end{swisstable}

Ein kurzer Absatz Grundtext zwischen den beiden Tabellen.

\begin{swisstable}{Zweite Legende, ebenfalls ueber mehrere Marginalzeilen gebrochen}
\begin{tabularx}{\textcolumn}[t]{@{}S{30mm}@{\hspace{4mm}}L@{}}
\toprule
\tabhead{Merkmal} & \tabhead{Beschreibung} \\
\midrule
Drittes & Ein weiterer Wert \\
Viertes & Und ein letzter \\
\bottomrule
\end{tabularx}
\end{swisstable}
\end{document}
